\chapter{Comenzi utile pentru reproducerea experimentelor}
\label{appendix:commands}

Această anexă listează comenzile principale folosite în proiect. Ele sunt utile pentru refacerea fluxului experimental pe orice sistem cu Windows PowerShell.

\section{Pornirea mediului}

\begin{lstlisting}[language=bash,caption={Inițializarea mediului Python},label={lst:venv}]
python -m venv .venv
.\.venv\Scripts\Activate.ps1
python -m pip install -e .
python -m playwright install chromium
\end{lstlisting}

\section{Colectarea datelor}

\begin{lstlisting}[language=bash,caption={Exemplu de scraping Autovit},label={lst:scrape-autovit}]
python -m carvaluator_scraper.cli scrape-search autovit "https://www.autovit.ro/autoturisme" --pages 150 --delay 0.4 --output data\autovit_xl.jsonl
\end{lstlisting}

\begin{lstlisting}[language=bash,caption={Exemplu de scraping mobile.de},label={lst:scrape-mobilede}]
python -m carvaluator_scraper.cli scrape-search mobilede "https://suchen.mobile.de/fahrzeuge/search.html?dam=false&isSearchRequest=true&ref=quickSearch&s=Car&vc=Car" --pages 5 --delay 1 --output data\mobilede_search.jsonl
\end{lstlisting}

\begin{lstlisting}[language=bash,caption={Scraping experimental mobile.de prin pagini Markdown},label={lst:scrape-mobilede-reader}]
python scripts\scrape_mobilede_reader.py --output data\mobilede_reader_20260606.jsonl --report data\mobilede_reader_20260606_report.json --cache-dir data\mobilede_reader_cache_20260606
\end{lstlisting}

\section{Pregătirea setului de date}

\begin{lstlisting}[language=bash,caption={Export CSV pentru antrenare},label={lst:export-csv}]
python -m carvaluator_scraper.cli export-csv data\autovit_xl.jsonl --output data\autovit_xl.csv --drop-fuzzy-duplicates --report data\autovit_xl_export_report.json
\end{lstlisting}

\begin{lstlisting}[language=bash,caption={Export CSV pentru setul mobile.de colectat prin reader},label={lst:export-mobilede-reader}]
python -m carvaluator_scraper.cli export-csv data\mobilede_reader_20260606.jsonl --output data\mobilede_reader_20260606.csv --drop-fuzzy-duplicates --report data\mobilede_reader_20260606_export_report.json
\end{lstlisting}

\begin{lstlisting}[language=bash,caption={Export CSV combinat pentru experimentul final},label={lst:export-combined-reader}]
python -m carvaluator_scraper.cli export-csv data\autovit_xl.jsonl data\autovit_more_20260526.jsonl data\mobilede_consolidated_20260606.jsonl data\mobilede_reader_20260606.jsonl --output data\combined_reader_20260606.csv --drop-fuzzy-duplicates --report data\combined_reader_20260606_report.json
\end{lstlisting}

\section{Antrenarea modelelor}

\begin{lstlisting}[language=bash,caption={Antrenare cu țintă logaritmică},label={lst:train-log}]
python -m carvaluator_scraper.cli train-models data\combined_reader_20260606.csv --output-dir data\model_results_combined_reader_20260606_log --log-target
\end{lstlisting}

\section{Pornirea aplicației}

\begin{lstlisting}[language=bash,caption={Pornire automată cu scriptul proiectului},label={lst:start-script}]
.\scripts\start-carvaluator.ps1 -Port 8017 -Background -StopExisting
\end{lstlisting}

\begin{lstlisting}[language=bash,caption={Pornirea serverului local},label={lst:start-server}]
$env:CARVALUATOR_MODEL_BUNDLE = "data\model_results_combined_reader_20260606_log\best_model.joblib"
$env:CARVALUATOR_SIMILARITY_CSV = "data\similarity_current.csv"
$env:CARVALUATOR_API_PORT = "8017"
.\.venv\Scripts\python.exe -m carvaluator_scraper.api
\end{lstlisting}

După pornire, aplicația rulează local la adresa \texttt{http://127.0.0.1:8017}. Ruta \texttt{/health} verifică dacă modelul, datasetul de similaritate și graficele principale sunt disponibile.

\section{Reîmprospătarea sugestiilor similare}

Anunțurile din baza de similaritate pot expira chiar dacă rândurile rămân valide pentru analiza statistică. Din acest motiv, lotul folosit de KNN trebuie actualizat periodic, separat de antrenarea modelelor. Comanda următoare colectează anunțuri \autovit\ curente, normalizează și deduplică rezultatul, arhivează fișierele cu timestamp, actualizează \texttt{data/similarity\_current.csv} și repornește serverul:

\begin{lstlisting}[language=bash,caption={Actualizarea locală a lotului de sugestii},label={lst:refresh-similarity-local}]
.\scripts\refresh-similarity.ps1 -Pages 100 -RestartServer
\end{lstlisting}

Actualizarea nu modifică estimarea de preț și nu reantrenează modelele. Ea schimbă numai mulțimea de candidați din care motorul KNN selectează anunțurile similare. Rapoartele deja salvate în istoricul conturilor păstrează sugestiile existente la momentul analizei.

Pentru ca același lot să fie inclus și în deployment-ul Render se folosește opțiunea suplimentară:

\begin{lstlisting}[language=bash,caption={Actualizarea sugestiilor locale și a artefactelor Render},label={lst:refresh-similarity-render}]
.\scripts\refresh-similarity.ps1 -Pages 100 -UpdateRenderArtifacts -RestartServer
\end{lstlisting}

Această comandă actualizează doar folderul local \texttt{deploy\_artifacts}. Site-ul public se schimbă abia după commit, push și redeploy, deoarece instanța Render nu poate citi fișierele din folderul local \texttt{data}.

\section{Pregătirea pentru Render}

\begin{lstlisting}[language=bash,caption={Pregătirea artefactelor pentru deployment Render},label={lst:render-artifacts}]
.\scripts\prepare-render-artifacts.ps1 -Clean
\end{lstlisting}

Comanda creează folderul \texttt{deploy\_artifacts}, copiază modelul \texttt{best\_model.joblib}, datasetul combinat și graficele generate la antrenare. Fișierul \texttt{render.yaml} folosește următoarele căi:

\begin{lstlisting}[language=bash,caption={Variabile principale pentru Render},label={lst:render-vars}]
CARVALUATOR_MODEL_BUNDLE=deploy_artifacts/model_results/best_model.joblib
CARVALUATOR_SIMILARITY_CSV=deploy_artifacts/datasets/similarity_current.csv
CARVALUATOR_TRUST_PROXY_HEADERS=1
CARVALUATOR_RATE_LIMIT_ENABLED=1
CARVALUATOR_RATE_LIMIT_GLOBAL_REQUESTS=60
CARVALUATOR_RATE_LIMIT_GLOBAL_WINDOW_SECONDS=60
CARVALUATOR_RATE_LIMIT_LOGIN_REQUESTS=10
CARVALUATOR_RATE_LIMIT_LOGIN_WINDOW_SECONDS=900
CARVALUATOR_RATE_LIMIT_REGISTER_REQUESTS=3
CARVALUATOR_RATE_LIMIT_REGISTER_WINDOW_SECONDS=3600
CARVALUATOR_RATE_LIMIT_PREDICT_REQUESTS=5
CARVALUATOR_RATE_LIMIT_PREDICT_WINDOW_SECONDS=60
\end{lstlisting}

Valorile de rate limiting corespund limitelor generale, de login, înregistrare și predicție. Variabila pentru antetele proxy trebuie activată numai în spatele proxy-ului Render; local poate rămâne dezactivată. Pentru deployment, scriptul produce un bundle redus de aproximativ 1 MB, care păstrează SVR, Ridge, KNN și Gradient Boosting. Modelul complet local, cu șapte pipeline-uri, are aproximativ 606 MB și nu este încărcat pe planul gratuit Render. Configurația curentă folosește Docker pentru a fixa versiunile scikit-learn și pentru a instala Chromium:

\begin{lstlisting}[language=bash,caption={Fișierele folosite de Blueprint-ul Render},label={lst:render-docker}]
Dockerfile
render.yaml
deploy_artifacts/model_results/best_model.joblib
deploy_artifacts/datasets/similarity_current.csv
\end{lstlisting}

După publicarea modificărilor în GitHub se folosește opțiunea de ștergere a cache-ului urmată de redeploy. Ruta \texttt{/health} trebuie să confirme existența modelului și a datasetului, precum și versiunile scikit-learn 1.8.0 și Playwright 1.58.0. Reîmprospătarea lotului reduce probabilitatea linkurilor expirate, dar nu o poate elimina complet, deoarece un anunț se poate închide imediat după colectare.
